% !TeX root = ../papers/nc-p3.tex
\section{正次元安定群を持つ有限候補の全次数正則性}
\label{sec:finite-connected-regularity}
\begin{theorem}\label{thm:finite-connected-regularity}
四辺の点対応が有限な標準有限Frobenius窓で、射影安定群が正次元なら、
全次数で標準quadratic Koszul左右AS正則であり、群作用を忘れた窓は半安定である。
加法群が作用すればそのJordan型は四位置で等しく、次のいずれかである。
\[
\begin{array}{c|c|l}
V_i\text{ の型}&R_i\text{ の型}&\text{多重安定極限}\\\hline
(4)&(5,1)&\text{多項式環}\\
(3,1)&(3,3)&\text{多項式環}\\
(2,1,1)&(2,2,1,1)&\text{多項式環または対称型の }P=0
\end{array}
\]
これは有限候補を入力とする判定で、各行の有限例の存在を必要としない。
乗法群の場合は定理\ref{thm:simple-weight-finite-windows}と
\ref{thm:double-middle-finite-windows}の既出の二曲線と一曲面に尽きる。
\end{theorem}
\begin{proof}[異なるJordan型の辺の排除]
前節の $D$ を使い、各Jordanブロックを $L_d$ に延長してその対角重みを使う。
各位置の最大ブロックの最低文字を $z$、最高文字を $x$ と書く。
$zz$ 項は必ず零である。$(3,1)$ と $(2,1,1)$ の順に隣接すると、
最低となり得る重み $-2$ の成分は $z\otimes U$（$U$ は第二位置の二自明文字）
だけである。$D^2r$ はそこから来る $x\otimes U$ と重み3の $kx\otimes x$ だけに入る。
純関係がないので $D^2 R=0$ となり、$z$ での評価が $U$ を含めない。
逆順も同じである。$(4)$ と $(3,1)$ の辺では最低 $-5$ が消え、
重み $-3$ の成分は $L_6$ の次最低と二個の $L_4$ の最低の直和にある。
$D^4R$ は最高 $kx\otimes x$ に入り零なので $L_6$ の成分はない。
$z$ での評価は第二位置の二つの重み0をともに含むから、独立な二つの最低成分が必要になる。
$D^j$, $0\le j<4$ によるそれらの像は異なる重み層ごとに独立なので $\dim R\ge8$ となって矛盾する。
$(4)$ と $(2,1,1)$ は命題\ref{prop:finite-additive-jordan}によっても排除済みである。

\end{proof}
\begin{proof}[同じJordan型の初期関係]
型 $(4)$ では $V_i\otimes V_{i+1}=L_7+L_5+L_3+L_1$ で、最低 $-6$ は消える。
$D^5 R\subset kx\otimes x$ は零なので、重み $-4$ の成分は $L_5$ の最低文字だけである。
正準点で次の文字への評価が必要なためこの成分は現れる。
従って $R$ は長さ5の鎖と長さ1の鎖である。四位置の同じ対角一径数群で窓 $w$ を還元する。
この極限は常に存在する。実際、加法群同変な各辺の埋込み
$L_5+L_1\to V_i\otimes V_{i+1}$ は重みを下げず、次数零写像は半単純群同変である
（前節の最高重みの議論）。ここで長さ5の始元は非零の重み $-4$ 成分を持つように取る。
相補辺間の抽象的な完全対合の次数零部分は同型で、特に重複度1の $L_5$ 間で非零である。
従って極限 $w_0$ は非零である。その全辺の関係像は $L_5+L_1=\bigwedge^2 L_4$ に入る。
これは四つの $L_4$ を同一視したとき各隣接交換で交代的ということで、$w_0$ は交代体積となる。
両切断で階数6が回復し、命題\ref{prop:unrestricted-good-reduction}が使える。

型 $(3,1)$ では重み $-4$ が消え、$D^3 R$ は最高の純直線に入り零である。
重み $-2$ 部分で $L_5$ の次最低が消え、残るのは三つの $L_3$ の最低である。
$z$ で評価すると相手の重み0の二文字へ全射なので独立な最低二成分がある。
これらとその一次、二次の $D$ 像だけで六次元を使い尽くす。
従って $R=2L_3$ と読んだ埋込みの次数零部分は単射で、両切断で $w_0$ は階数6を保つ。
その重複度二平面は、三コピー $t_i X,Xt_{i+1},X\times X$ のうち、
最後と最初又は最後と二番目を取る両射影で同型になる。
これは元の辺の正準点で左右それぞれ、$t$ と中間の鎖文字の両方への評価が全射であるためである。
特にこの二平面には最初又は二番目だけのベクトルはなく、各文字の支持も全体である。
定理\ref{thm:sl2-window-profiles}の $L_3+L_1$ の証明で、階数2かつ全一文字支持でも
非正則になり得るのは $uv\ne0$ である $a_i$ が消える場合に限られる。このとき隣辺に純混合コピー
$t_i X$ 又は $X t_i$ が入るが、上の二射影がこれを排除する。
従って $w_0$ は多項式環または $U(\mathfrak{sl}_2)$ の斉次化である。
良い還元と系\ref{cor:semisimple-symmetry-closed-limits}で正則性と閉極限を得る。
型 $(2,1,1)$ は定理\ref{thm:finite-four-transvections}である。
\end{proof}

\begin{proof}[正則性を仮定しない乗法群の重み]
定理\ref{thm:finite-torus-profiles}の露出面、端点評価、辺ごとの局所指標の証明は
有限完全交叉と純関係の排除だけによる。そこで正則窓のEuler文字を用いた箇所を次の二点で置き換える。
第一に、最低が三重なら隣接位置でも三重であり、最低の $3\times3$ 部分の関係は
純テンソルを避ける上界4と有限交叉の下界4から四次元である。
全六関係中、最高関係重みの重複度は高々2となる。反対辺と窓文字への完全対合では
最低と最高の関係重複度が対応するため不可能である。最高が三重の場合も反転すればよい。

第二に、両極単純で中間も単純なとき、辺の指標の列挙は中央差を $1,2$ に制限し、
差2の隣接を禁じる。異なる差の辺では関係重み $1,\ldots,6$ が各一本で、両位置の
両極どうし以外の組は辺の点ではない。一方、差1どうしの辺の中央関係重複度は2である。
よって混合辺が一つでもあれば、その反対辺も混合辺である。四位置で差2は交互に並ぶしかない。
このとき第一切断で「最高、最低、最低、最高」の係数を見る。行の位置対も列の位置対も
グラフにないから、それぞれの関係線のその座標は非零である。
その二重みは補い合う。実際全体重みは $1+6=7$ で、二つの最高文字は重み3と4である。
従ってこの係数は階数1の完全対合により非零となる。他方、次の切断では中央の最低どうしに
関係が一つもなく、この係数は零である。矛盾する。
中央二重の議論は同定理の有限交叉のままである。残る二型が満たす係数条件とPBW性は
前掲の二分類定理の切断計算による。特に中央二重の両極の交互語係数は、二つの極の包の積に
中央重みの関係が二本独立に制限されることと、同じ極の連続を含む語の係数が零であることから非零である。
最後に標数零の正次元線形代数群は乗法群又は加法群を含むので、以上ですべてを覆う。
\end{proof}
